\chapter{Tablas completas de resultados}
\label{ap:tablas}

Se recogen aquí los resultados desagregados por conjunto de datos. Todas las
cifras provienen de la misma ejecución del banco de pruebas descrito en el
\cref{cap:resultados} y son reproducibles con el procedimiento del
\cref{ap:software}.

Los cuatro conjuntos son: los \datoNSinteticos{} escenarios sintéticos con la
tolerancia de producción; las \datoNReales{} pólizas reales de renta vitalicia;
los \datoNRealesRP{} casos procedentes del conjunto de retiro programado,
evaluados en el catálogo con la misma función de valoración de renta vitalicia;
y de nuevo los escenarios sintéticos con tolerancia estricta
($\varepsilon_f=10^{-6}$), que permite contrastar el trabajo requerido al
exigir un residuo menor. Las sucesiones de pocos pasos o con terminación
finita no permiten estimar un orden asintótico.

\begin{center}
\small
\begin{longtable}{@{}llrrrrr@{}}
\toprule
\textbf{Conjunto} & \textbf{Método} & \textbf{Éxito} & \textbf{Ev.} & \textbf{Iter.} & \textbf{ms} & \textbf{Error mediano} \\
\midrule
\endhead
Pólizas RV & Newton semisuave & 100\% & 2{,}78 & 1{,}78 & 14{,}2 & $4{,}7\cdot10^{-10}$ \\
Pólizas RV & Punto fijo relajado & 100\% & 2{,}78 & 1{,}78 & 14{,}3 & $4{,}7\cdot10^{-10}$ \\
Pólizas RV & Newton (analítico) & 100\% & 2{,}78 & 1{,}78 & 14{,}7 & $4{,}7\cdot10^{-10}$ \\
Pólizas RV & Steffensen & 0\% & 3{,}00 & 1{,}00 & 16{,}0 & $1{,}4\cdot10^{6}$ \\
Pólizas RV & Aitken--Steffensen & 100\% & 4{,}22 & 1{,}78 & 21{,}1 & $4{,}7\cdot10^{-10}$ \\
Pólizas RV & Steffensen escalado & 100\% & 4{,}56 & 1{,}78 & 25{,}1 & $4{,}7\cdot10^{-10}$ \\
Pólizas RV & Newton salvaguardado & 100\% & 4{,}78 & 1{,}78 & 25{,}1 & $4{,}7\cdot10^{-10}$ \\
Pólizas RV & Anderson--Björck & 100\% & 4{,}78 & 2{,}78 & 26{,}3 & $2{,}9\cdot10^{-6}$ \\
Pólizas RV & Pegasus & 100\% & 4{,}78 & 2{,}78 & 25{,}1 & $4{,}6\cdot10^{-6}$ \\
Pólizas RV & Secante & 100\% & 4{,}78 & 3{,}78 & 23{,}6 & $4{,}7\cdot10^{-10}$ \\
Pólizas RV & Brent & 100\% & 5{,}00 & 3{,}00 & 25{,}0 & $0$ \\
Pólizas RV & Regula falsi & 100\% & 5{,}33 & 3{,}33 & 26{,}3 & $3{,}7\cdot10^{-6}$ \\
Pólizas RV & Illinois & 100\% & 5{,}56 & 3{,}56 & 29{,}0 & $4{,}7\cdot10^{-10}$ \\
Pólizas RV & Müller & 100\% & 5{,}78 & 4{,}78 & 31{,}3 & $8{,}4\cdot10^{-6}$ \\
Pólizas RV & Newton (dif.\ finitas) & 100\% & 7{,}00 & 2{,}00 & 35{,}7 & $3{,}6\cdot10^{-6}$ \\
Pólizas RV & Ridders & 100\% & 9{,}11 & 3{,}56 & 53{,}2 & $4{,}7\cdot10^{-10}$ \\
Pólizas RV & Halley & 100\% & 10{,}44 & 3{,}22 & 51{,}9 & $4{,}9\cdot10^{-5}$ \\
Pólizas RV & ITP & 100\% & 29{,}56 & 27{,}56 & 148{,}1 & $1{,}6\cdot10^{-3}$ \\
Pólizas RV & Bisección & 100\% & 32{,}67 & 30{,}67 & 162{,}8 & $1{,}6\cdot10^{-3}$ \\
\midrule
Registros adicionales & Newton semisuave & 100\% & 2{,}71 & 1{,}71 & 14{,}5 & $9{,}3\cdot10^{-10}$ \\
Registros adicionales & Punto fijo relajado & 100\% & 2{,}71 & 1{,}71 & 13{,}9 & $9{,}3\cdot10^{-10}$ \\
Registros adicionales & Newton (analítico) & 100\% & 2{,}71 & 1{,}71 & 14{,}6 & $9{,}3\cdot10^{-10}$ \\
Registros adicionales & Steffensen & 0\% & 3{,}00 & 1{,}00 & 15{,}2 & $1{,}9\cdot10^{6}$ \\
Registros adicionales & Aitken--Steffensen & 100\% & 4{,}00 & 1{,}71 & 20{,}8 & $9{,}3\cdot10^{-10}$ \\
Registros adicionales & Steffensen escalado & 100\% & 4{,}43 & 1{,}71 & 23{,}1 & $1{,}4\cdot10^{-9}$ \\
Registros adicionales & Newton salvaguardado & 100\% & 4{,}71 & 1{,}71 & 25{,}5 & $9{,}3\cdot10^{-10}$ \\
Registros adicionales & Anderson--Björck & 100\% & 4{,}71 & 2{,}71 & 25{,}7 & $4{,}7\cdot10^{-6}$ \\
Registros adicionales & Pegasus & 100\% & 4{,}71 & 2{,}71 & 26{,}3 & $6{,}6\cdot10^{-6}$ \\
Registros adicionales & Secante & 100\% & 4{,}71 & 3{,}71 & 27{,}1 & $4{,}7\cdot10^{-10}$ \\
Registros adicionales & Brent & 100\% & 5{,}00 & 3{,}00 & 25{,}5 & $0$ \\
Registros adicionales & Illinois & 100\% & 5{,}43 & 3{,}43 & 27{,}2 & $4{,}7\cdot10^{-10}$ \\
Registros adicionales & Regula falsi & 100\% & 5{,}43 & 3{,}43 & 29{,}0 & $1{,}9\cdot10^{-6}$ \\
Registros adicionales & Müller & 100\% & 5{,}71 & 4{,}71 & 30{,}7 & $1{,}3\cdot10^{-5}$ \\
Registros adicionales & Newton (dif.\ finitas) & 100\% & 6{,}57 & 1{,}86 & 36{,}4 & $3{,}8\cdot10^{-6}$ \\
Registros adicionales & Ridders & 100\% & 8{,}86 & 3{,}43 & 47{,}3 & $4{,}7\cdot10^{-10}$ \\
Registros adicionales & Halley & 100\% & 9{,}57 & 2{,}86 & 53{,}3 & $4{,}3\cdot10^{-5}$ \\
Registros adicionales & ITP & 100\% & 29{,}71 & 27{,}71 & 166{,}6 & $6{,}0\cdot10^{-4}$ \\
Registros adicionales & Bisección & 100\% & 34{,}00 & 32{,}00 & 178{,}2 & $1{,}7\cdot10^{-3}$ \\
\midrule
Sintéticos & Steffensen & 0\% & 3{,}00 & 1{,}00 & 17{,}7 & $2{,}5\cdot10^{6}$ \\
Sintéticos & Newton semisuave & 100\% & 3{,}67 & 2{,}67 & 20{,}2 & $0$ \\
Sintéticos & Newton (analítico) & 100\% & 3{,}67 & 2{,}67 & 22{,}2 & $0$ \\
Sintéticos & Aitken--Steffensen & 100\% & 4{,}67 & 2{,}00 & 25{,}8 & $0$ \\
Sintéticos & Punto fijo relajado & 100\% & 4{,}78 & 3{,}78 & 25{,}8 & $3{,}8\cdot10^{-6}$ \\
Sintéticos & Pegasus & 100\% & 5{,}67 & 3{,}67 & 30{,}7 & $3{,}6\cdot10^{-5}$ \\
Sintéticos & Secante & 100\% & 5{,}67 & 4{,}67 & 31{,}4 & $9{,}3\cdot10^{-10}$ \\
Sintéticos & Anderson--Björck & 100\% & 5{,}67 & 3{,}67 & 30{,}5 & $3{,}2\cdot10^{-5}$ \\
Sintéticos & Brent & 100\% & 5{,}78 & 3{,}78 & 30{,}8 & $2{,}4\cdot10^{-7}$ \\
Sintéticos & Illinois & 100\% & 6{,}33 & 4{,}33 & 34{,}1 & $2{,}3\cdot10^{-10}$ \\
Sintéticos & Regula falsi & 100\% & 6{,}56 & 4{,}56 & 35{,}3 & $3{,}2\cdot10^{-4}$ \\
Sintéticos & Steffensen escalado & 100\% & 6{,}56 & 2{,}78 & 35{,}4 & $4{,}7\cdot10^{-10}$ \\
Sintéticos & Müller & 100\% & 6{,}67 & 5{,}67 & 37{,}1 & $1{,}0\cdot10^{-5}$ \\
Sintéticos & Newton salvaguardado & 100\% & 6{,}89 & 3{,}89 & 38{,}0 & $0$ \\
Sintéticos & Newton (dif.\ finitas) & 100\% & 9{,}00 & 2{,}67 & 54{,}6 & $2{,}2\cdot10^{-8}$ \\
Sintéticos & Halley & 100\% & 10{,}33 & 3{,}11 & 55{,}8 & $9{,}3\cdot10^{-10}$ \\
Sintéticos & Ridders & 100\% & 11{,}33 & 4{,}67 & 59{,}7 & $4{,}2\cdot10^{-6}$ \\
Sintéticos & ITP & 100\% & 29{,}67 & 27{,}67 & 164{,}3 & $1{,}4\cdot10^{-3}$ \\
Sintéticos & Bisección & 100\% & 33{,}67 & 31{,}67 & 199{,}4 & $1{,}4\cdot10^{-3}$ \\
\midrule
Sint.\ tol.\ fina & Steffensen & 0\% & 3{,}00 & 1{,}00 & 17{,}2 & $2{,}5\cdot10^{6}$ \\
Sint.\ tol.\ fina & Newton semisuave & 100\% & 3{,}67 & 2{,}67 & 22{,}0 & $0$ \\
Sint.\ tol.\ fina & Newton (analítico) & 100\% & 3{,}67 & 2{,}67 & 23{,}2 & $0$ \\
Sint.\ tol.\ fina & Aitken--Steffensen & 100\% & 4{,}67 & 2{,}00 & 26{,}8 & $0$ \\
Sint.\ tol.\ fina & Secante & 100\% & 6{,}00 & 5{,}00 & 33{,}7 & $9{,}3\cdot10^{-10}$ \\
Sint.\ tol.\ fina & Brent & 100\% & 6{,}44 & 4{,}44 & 38{,}7 & $0$ \\
Sint.\ tol.\ fina & Steffensen escalado & 100\% & 6{,}56 & 2{,}78 & 38{,}8 & $4{,}7\cdot10^{-10}$ \\
Sint.\ tol.\ fina & Anderson--Björck & 100\% & 6{,}67 & 4{,}67 & 39{,}9 & $9{,}3\cdot10^{-10}$ \\
Sint.\ tol.\ fina & Pegasus & 100\% & 6{,}67 & 4{,}67 & 37{,}6 & $4{,}7\cdot10^{-10}$ \\
Sint.\ tol.\ fina & Punto fijo relajado & 100\% & 6{,}78 & 5{,}78 & 37{,}5 & $9{,}3\cdot10^{-10}$ \\
Sint.\ tol.\ fina & Illinois & 100\% & 6{,}78 & 4{,}78 & 39{,}1 & $0$ \\
Sint.\ tol.\ fina & Newton salvaguardado & 100\% & 6{,}89 & 3{,}89 & 44{,}2 & $0$ \\
Sint.\ tol.\ fina & Müller & 100\% & 7{,}67 & 6{,}67 & 44{,}5 & $9{,}3\cdot10^{-10}$ \\
Sint.\ tol.\ fina & Regula falsi & 100\% & 8{,}78 & 6{,}78 & 49{,}3 & $1{,}9\cdot10^{-9}$ \\
Sint.\ tol.\ fina & Newton (dif.\ finitas) & 100\% & 10{,}67 & 3{,}22 & 64{,}4 & $4{,}7\cdot10^{-10}$ \\
Sint.\ tol.\ fina & Halley & 100\% & 11{,}33 & 3{,}44 & 70{,}0 & $0$ \\
Sint.\ tol.\ fina & Ridders & 100\% & 14{,}00 & 6{,}00 & 77{,}0 & $9{,}3\cdot10^{-10}$ \\
Sint.\ tol.\ fina & ITP & 100\% & 32{,}67 & 30{,}67 & 192{,}9 & $9{,}3\cdot10^{-10}$ \\
Sint.\ tol.\ fina & Bisección & 100\% & 51{,}22 & 49{,}22 & 328{,}3 & $9{,}3\cdot10^{-10}$ \\
\midrule
\bottomrule
\end{longtable}
\end{center}

\section*{Lectura del conjunto de tolerancia estricta}

El contraste entre el primer bloque (tolerancia de producción,
$\varepsilon_f=\datoEpsilonProd$) y el último (tolerancia estricta,
$\varepsilon_f=10^{-6}$) permite examinar la sensibilidad del costo a la
precisión solicitada:

\begin{itemize}[leftmargin=*]
  \item Para la bisección, reducir por diez la tolerancia de error certificada
        por el intervalo exige unas $\log_2(10)\approx3{,}3$ reducciones
        adicionales en su cota de peor caso (\cref{thm:cota-biseccion}). Una
        parada anterior por residuo puede reducir ese número. Esta cota no
        clasifica a ITP como método de orden lineal ni describe por sí sola
        el costo de regula falsi.
  \item Para los métodos que aprovechan el tramo afín de la raíz, el
        \cref{thm:terminacion-finita} explica que una tolerancia más estricta
        pueda requerir pocos pasos adicionales o ninguno en aritmética exacta.
        En el cómputo efectivo intervienen el redondeo y las evaluaciones
        necesarias para certificar la parada; los conteos de la tabla son
        observaciones de estos casos, no una estimación de orden asintótico.
\end{itemize}
